% =============================================================
% Gemeinsamer Buch-Rumpf für alle Master-Dateien.
%
% Die Master-Dateien (geheimnis-der-redundanz.tex und alle
% redundanz-tagN.tex) setzen die zwei Steuer-Makros …
%
%   \def\maxetappe{N}              % wieviele Etappen ins Buch?
%   \def\hauptbuchtitelseite{...}  % Titelseite-Block, leer für Standalones
%
% … laden danach den scrbook-Klassenkopf (das passiert in jedem
% Master direkt nach \documentclass) und greifen schließlich per
% % =============================================================
% Gemeinsamer Buch-Rumpf für alle Master-Dateien.
%
% Die Master-Dateien (geheimnis-der-redundanz.tex und alle
% redundanz-tagN.tex) setzen die zwei Steuer-Makros …
%
%   \def\maxetappe{N}              % wieviele Etappen ins Buch?
%   \def\hauptbuchtitelseite{...}  % Titelseite-Block, leer für Standalones
%
% … laden danach den scrbook-Klassenkopf (das passiert in jedem
% Master direkt nach \documentclass) und greifen schließlich per
% % =============================================================
% Gemeinsamer Buch-Rumpf für alle Master-Dateien.
%
% Die Master-Dateien (geheimnis-der-redundanz.tex und alle
% redundanz-tagN.tex) setzen die zwei Steuer-Makros …
%
%   \def\maxetappe{N}              % wieviele Etappen ins Buch?
%   \def\hauptbuchtitelseite{...}  % Titelseite-Block, leer für Standalones
%
% … laden danach den scrbook-Klassenkopf (das passiert in jedem
% Master direkt nach \documentclass) und greifen schließlich per
% % =============================================================
% Gemeinsamer Buch-Rumpf für alle Master-Dateien.
%
% Die Master-Dateien (geheimnis-der-redundanz.tex und alle
% redundanz-tagN.tex) setzen die zwei Steuer-Makros …
%
%   \def\maxetappe{N}              % wieviele Etappen ins Buch?
%   \def\hauptbuchtitelseite{...}  % Titelseite-Block, leer für Standalones
%
% … laden danach den scrbook-Klassenkopf (das passiert in jedem
% Master direkt nach \documentclass) und greifen schließlich per
% \input{buch-rumpf} auf den hier definierten Body zu.
% =============================================================

\input{praeambel}
\input{farben}
\input{befehle}
\input{glossar}

% Defaults, falls ein Master vergisst sie zu setzen.
\providecommand{\maxetappe}{8}
\providecommand{\hauptbuchtitelseite}{}

\begin{document}

\frontmatter
\hauptbuchtitelseite
\pagestyle{plain}

\include{vorwort}
\tableofcontents

\mainmatter
\pagestyle{scrheadings}

% Etappen 1 .. \maxetappe der Reihe nach einbinden. \include sorgt
% pro Datei für eigenen aux-File und Pagebreak.
\foreach \etappennr in {1,...,\maxetappe}{%
  \include{kapitel/tag\etappennr}%
}

\appendix

% Kapitel-spezifische Anhänge (z. B. Tabellen, Werkstatt-Material).
% Pro Etappe wird anhang/tagN.tex eingebunden, wenn die Datei existiert.
\foreach \etappennr in {1,...,\maxetappe}{%
  \IfFileExists{anhang/tag\etappennr.tex}{%
    \input{anhang/tag\etappennr}%
  }{}%
}

% Lösungen zu allen bisher behandelten Etappen, gesammelt im Anhang.
\chapter{Lösungen}

\begin{quote}
Erst nach eigenem Versuch ansehen — der Aha-Moment beim Selber-Lösen
ist wichtiger als die fertige Lösung am Ende.
\end{quote}

\foreach \etappennr in {1,...,\maxetappe}{%
  \input{loesungen/tag\etappennr}%
}

% \glsaddall zeigt alle Glossar-Einträge (auch die, die im Text nicht
% per \gls{...} referenziert wurden). Solange wir die Markups nicht
% systematisch durch alle Kapitel gezogen haben, ist das die
% praktikable Variante.
\glsaddall
\printglossary[title=Glossar, toctitle=Glossar]

\printindex

\end{document}
 auf den hier definierten Body zu.
% =============================================================

% =============================================================
% Präambel — Pakete, Schriften, Geometrie, Sprache, Kopf-/Fußzeile.
%
% Wird von buch.tex eingebunden, NACH \documentclass.
% Layout-Bausteine (Aufgaben-Umgebungen) sind in befehle.tex,
% Farben in farben.tex.
% =============================================================

% --- Geometrie nach doku/LAYOUT_BRIEFING.md ---
\usepackage[a4paper,
            inner=2.5cm, outer=4.5cm,
            top=2.0cm, bottom=2.5cm,
            marginparwidth=2.7cm, marginparsep=0.15cm,
            headheight=15pt]{geometry}

% --- Schriften ---
% Source Pro Familie als Hausschriften, vendored unter latex/fonts/
% (siehe latex/fonts/README.md für Provenance). Source Serif für den
% Fließtext, Source Sans für Überschriften/Marginalien, Source Code
% für Quelltexte. Math-Formeln werden weiter unten via unicode-math
% gesetzt (STIX Two Math passt visuell zu Source Serif).
\usepackage{fontspec}
\setmainfont{SourceSerif4}[
  Path=fonts/source-serif/,
  Extension=.otf,
  UprightFont=*-Regular,
  ItalicFont=*-It,
  BoldFont=*-Bold,
  BoldItalicFont=*-BoldIt,
  Scale=1.00,
]
\setsansfont{SourceSans3}[
  Path=fonts/source-sans/,
  Extension=.otf,
  UprightFont=*-Regular,
  ItalicFont=*-It,
  BoldFont=*-Bold,
  BoldItalicFont=*-BoldIt,
  Scale=0.95,
]
\setmonofont{SourceCodePro}[
  Path=fonts/source-code-pro/,
  Extension=.otf,
  UprightFont=*-Regular,
  ItalicFont=*-It,
  BoldFont=*-Bold,
  BoldItalicFont=*-BoldIt,
  Scale=0.88,
]
% Marginalien-Symbole (✏ ⌨ ⚙) brauchen einen Symbol-fähigen Font.
% Source Sans hat die Glyphen nicht; wir behalten DejaVu Sans als
% Fallback nur dafür (System-Font, im Container via fonts-dejavu).
\newfontfamily\markerfont{DejaVu Sans}

% --- Sprache, Mikrotypografie ---
\usepackage[ngerman]{babel}
% babel-ngerman macht ASCII-" zum Shorthand-Aktivator ("u → ü, "f → ff
% etc.). Das hat schon zu absurden Verschmelzungen wie "echterFFehler"
% statt "echter" Fehler" geführt. Wir schalten den Shorthand ab; für
% typografisch korrekte deutsche Anführungszeichen nutzen wir Unicode
% („…") direkt im Quelltext.
\shorthandoff{"}
\usepackage{microtype}

% --- Mathematik (\bmod, \le, \ge, \checkmark, …) ---
% amsmath für Mathematik-Umgebungen (align, gather, …).
% unicode-math + STIX Two Math ersetzt die Computer-Modern-Math-
% Schriften; Stilfamilie passt visuell zu Source Serif.
\usepackage{amsmath}
\usepackage{unicode-math}
\setmathfont{STIXTwoMath-Regular.otf}[
  Path=fonts/stix-two/,
]

% --- Glossar (Polish-Sprint M6) ---
% automake=immediate ruft makeglossaries per \write18 auf (shell-escape
% ist eh schon aktiv wegen minted), nonumberlist unterdrückt die
% Seitenzahlen hinter den Glossar-Einträgen (sonst stehen dort die
% Seiten, an denen \gls genutzt wurde).
\usepackage[automake=immediate, toc, nonumberlist]{glossaries-extra}
\makeglossaries

% --- Sachindex ---
\usepackage{imakeidx}
\makeindex[intoc, columns=2]

% --- Farben & Hyperlinks ---
\usepackage{xcolor}
\usepackage[hidelinks]{hyperref}
% pdfpagelabels ist hyperref-Default; wir setzen nur plainpages=false,
% damit hyperref beim \pagenumbering-Wechsel Frontmatter→Mainmatter
% keine Duplikat-Warnung wirft und PDF-Reader die Buchseiten-Labels
% (i, ii, …, 1, 2, …) korrekt anzeigen.
\hypersetup{plainpages=false}

% --- Layout-Bausteine ---
% tcolorbox: nur die Libraries laden, die wir tatsächlich brauchen
% (skins für `enhanced` und `borderline`, breakable für mehrseitige Boxen).
% [most] würde tikzfill u.a. mitziehen, was wir nicht nutzen.
\usepackage{tcolorbox}
\tcbuselibrary{skins, breakable}
\usepackage{changepage}
% \strictpagecheck macht \checkoddpage zuverlässig auch bei Page-Breaks
% (sonst kann die Seitenparität pro Page-Part falsch geprüft werden).
\strictpagecheck

% --- TikZ (für Venn-Diagramme etc.) ---
\usepackage{tikz}
\usetikzlibrary{positioning, calc}

% --- Tabellen ---
\usepackage{booktabs, array, longtable}

% --- Listen mit (a)/(b)/(c)-Labels für Aufgabenteile ---
\usepackage{enumitem}

% --- needspace: erzwingt Mindestplatz vor Aufgabenblöcken ---
\usepackage{needspace}

% --- Code-Highlighting via minted ---
% minted ruft Pygments per Shell auf; deshalb -shell-escape in .latexmkrc.
\usepackage{minted}
% Hintergrundfarbe für Code-Blöcke (sehr helles Grau, naive CMYK aus
% RGB(246, 247, 248), entspricht der bisherigen pandoc-shadecolor).
\definecolor{codebg}{cmyk}{0.01, 0.00, 0.00, 0.03}
\setminted{
  style=tango,
  fontsize=\small,
  bgcolor=codebg,
  frame=none,
  breaklines=true,
  autogobble=true,
  linenos=false,
  baselinestretch=1.05,
  xleftmargin=2pt,
  xrightmargin=2pt,
}

% --- Kopf-/Fußzeile via scrlayer-scrpage ---
\usepackage[autooneside=false, automark]{scrlayer-scrpage}
\clearpairofpagestyles
\automark[chapter]{chapter}
\ihead{\small\textit{\headmark}}
\ohead{\small\pagemark}
\KOMAoptions{headsepline=0.3pt}

% --- Witwen / Waisen vermeiden ---
\widowpenalty=10000
\clubpenalty=10000

% --- Etwas Toleranz für die schmaleren Spalten ---
\setlength{\emergencystretch}{3em}

% =============================================================
% Farbset 3 aus doku/LAYOUT_BRIEFING.md (Pastell, hell).
%
% Definiert in CMYK, weil das Buch in den Druck soll.
% Die ursprünglichen RGB-Werte aus dem Briefing stehen als Kommentar
% darüber; die CMYK-Werte sind naiv per Standardformel umgerechnet
% (K = 1 - max(R,G,B)/255 etc.). Für endgültige Druckqualität ggf.
% mit ICC-Profil (z. B. ISOcoated_v2_eci) nachjustieren.
%
% Diese drei Farben werden konsequent verwendet:
%   - für den Strich am Aufgabenblock
%   - für das Symbol in der Marginalie
%   - für die Kapitälchen-Bezeichnung unter dem Symbol
% =============================================================

% Bleistift — RGB(235, 160,  80), helles Orange
\definecolor{cBleistift}{cmyk}{0.00, 0.32, 0.66, 0.08}

% Python — RGB( 95, 140, 180), sanftes Blau
\definecolor{cPython}   {cmyk}{0.47, 0.22, 0.00, 0.29}

% Werkzeug — RGB(185, 200, 110), helles Lindgrün
\definecolor{cWerkzeug} {cmyk}{0.08, 0.00, 0.45, 0.22}

% =============================================================
% Aufgaben-Umgebungen.
%
% Drei Komfort-Umgebungen für die drei Aufgabentypen:
%   \begin{bleistiftuebung}{N}{Titel}  ...  \end{bleistiftuebung}
%   \begin{pythoneinheit} {N}{Titel}  ...  \end{pythoneinheit}
%   \begin{werkzeugcheck}    {Titel}  ...  \end{werkzeugcheck}
%
% Jede Umgebung erzeugt:
%   - eine farbige tcolorbox um den Aufgabentext, mit einem 2,5-pt-Strich
%     am AUSSEN-Rand der Seite (recto = rechts, verso = links),
%   - eine Marginalie auf derselben Außenseite mit Symbol und
%     Kapitälchen-Bezeichnung in der Aufgaben-Farbe.
%
% Die Bausteine stammen aus doku/LAYOUT_BRIEFING.md, dort steht auch
% die Begründung für jede Designentscheidung.
% =============================================================

% --- aufgabebox: tcolorbox mit page-aware Strich am Außenrand ---
% Eine einzige Box, kein Recto/Verso-Split mehr. Der overlay-Hook wird
% bei jedem Page-Part neu ausgewertet (auch wenn die Aufgabe per
% breakable über Seiten geht), prüft die aktuelle Seitenparität und
% zeichnet den Strich auf der jeweiligen Außenseite.
%
% left/right=8pt sorgt für Innenabstand zum Strich auf beiden Seiten;
% der Text bleibt immer mittig in der Box.
\newtcolorbox{aufgabebox}[1]{%
  enhanced, breakable,
  colback=white, colframe=white,
  boxrule=0pt, arc=0pt, outer arc=0pt,
  boxsep=0pt,
  left=8pt, right=8pt, top=4pt, bottom=4pt,
  before skip=0.6em, after skip=0.6em,
  overlay={%
    \checkoddpage
    \ifoddpage
      \draw[line width=2.5pt, color=#1]
        (frame.north east) -- (frame.south east);%
    \else
      \draw[line width=2.5pt, color=#1]
        (frame.north west) -- (frame.south west);%
    \fi
  },
}

% --- Marginalien-Inhalt für Recto und Verso ---
% Auf Recto-Seiten (Marginalspalte rechts): linksbündig zum Strich.
% Auf Verso-Seiten (Marginalspalte links):  rechtsbündig zum Strich.
% Das \vspace*{12pt} kompensiert das `before skip` und das `top` der
% tcolorbox, damit das Symbol mit der ersten Zeile der Aufgabenbox fluchtet.
\newcommand{\markercontentRecto}[4]{%
  \vspace*{12pt}\raggedright\color{#4}%
  {\markerfont\fontsize{24}{26}\selectfont #1}\\[0.05em]
  {\small\textsc{#2}}\\
  {\small\textsc{#3}}%
}
\newcommand{\markercontentVerso}[4]{%
  \vspace*{12pt}\raggedleft\color{#4}%
  {\markerfont\fontsize{24}{26}\selectfont #1}\\[0.05em]
  {\small\textsc{#2}}\\
  {\small\textsc{#3}}%
}

% --- Aufgaben-Umgebung ---
% Marginpar wird mit beiden Inhaltsvarianten aufgerufen (recto/verso);
% LaTeX wählt selbst die passende, je nachdem auf welcher Seite die
% Aufgabe beginnt. Der page-aware Strich entsteht im Overlay der
% aufgabebox und wechselt automatisch beim Seitenumbruch mit.
\newenvironment{aufgabe}[4]%
  {% #1 = Symbol, #2 = Label-Zeile-1, #3 = Label-Zeile-2, #4 = Farbe
   \marginpar[\markercontentVerso{#1}{#2}{#3}{#4}]%
              {\markercontentRecto{#1}{#2}{#3}{#4}}%
   \begin{aufgabebox}{#4}\ignorespaces}
  {\end{aufgabebox}}

% --- Komfort-Wrapper für die drei Aufgabentypen ---
% Konvention: der erste Teil der Aufgabe ist der Titel (fett), gefolgt
% vom Aufgabentext. Der Titel wird hier als Argument übergeben.
%
% Wir nutzen {\bfseries ...\par} statt \textbf{...\par}, weil \textbf
% als Inline-Befehl keinen Absatzumbruch (\par) im Argument verträgt
% ("Paragraph ended before \text@command was complete").
\newenvironment{bleistiftuebung}[2]%
  {\begin{aufgabe}{✏}{Bleistift-}{Übung #1}{cBleistift}{\bfseries #2.\par}}
  {\end{aufgabe}}

\newenvironment{pythoneinheit}[2]%
  {\begin{aufgabe}{⌨}{Python-}{Einheit #1}{cPython}{\bfseries #2.\par}}
  {\end{aufgabe}}

\newenvironment{werkzeugcheck}[1]%
  {\begin{aufgabe}{⚙}{Werkzeug-}{Check}{cWerkzeug}{\bfseries #1.\par}}
  {\end{aufgabe}}

% --- Code-Snippet-Einbindung ---
% \pythoncode{tagN/file.py} bindet eine Python-Datei vollständig ein.
% Pfad ist relativ zu latex/code/. Die .py-Dateien sind eigenständig
% lauffähig, sodass Greta sie 1:1 in Thonny verwenden kann.
\newcommand{\pythoncode}[1]{\inputminted{python}{code/#1}}

% \pythoncodeteil{tagN/file.py}{regionname} zeigt nur eine Region
% aus einer Datei. Region-Marker in der Quelle:
%     # region <name>
%     ... Code ...
%     # endregion
% Greta tippt die Datei sukzessive auf, im Buch zeigen wir pro
% Python-Einheit jeweils nur das aktuell relevante Stück.
%
% Implementierung: extract_region.py schneidet die Region zur
% Build-Zeit in eine temporäre Datei, die minted dann einbindet.
% Der Counter sorgt dafür, dass jede Einbindung ihre eigene Datei
% bekommt — nötig, weil minted das Snippet zur Build-Zeit cached.
\newcounter{snippetcntr}
\makeatletter
\newcommand{\pythoncodeteil}[2]{%
  \stepcounter{snippetcntr}%
  \protected@edef\@snippetfile{.snippets/snippet-\arabic{snippetcntr}.py}%
  \immediate\write18{python3 scripts/extract_region.py code/#1 #2 > \@snippetfile}%
  \inputminted{python}{\@snippetfile}%
}
\makeatother

% =============================================================
% Glossar-Einträge.
%
% Eingebunden in der Präambel (vor \begin{document}). Im Text
% verwendet via \gls{schluessel} an der ersten Erwähnung; weitere
% Auftritte können auch \gls{} sein, der Glossar-Eintrag wird dann
% nur noch normal gesetzt.
%
% Schlüssel sind kurz und sprechend; Anzeigename steht in `name=`.
% =============================================================

\newglossaryentry{modulo}{
  name=Modulo,
  description={Operation \(a \bmod n\): der Rest, der bei der Division
    von \(a\) durch \(n\) bleibt. Universalwerkzeug der
    Codierungstheorie}}

\newglossaryentry{paritaet}{
  name=Parität,
  description={Eigenschaft einer Bitfolge, die sich aus der Anzahl der
    Einsen ergibt: \emph{gerade Parität}, wenn die Anzahl gerade
    ist, sonst \emph{ungerade}. Eine Paritätsprüfung erkennt jede
    ungerade Anzahl gekippter Bits}}

\newglossaryentry{paritaetsbit}{
  name=Paritätsbit,
  description={Zusätzliches Bit, das so gewählt wird, dass die
    Gesamt-Parität (meist gerade) zustande kommt. Erkennt
    Einzel-Bitfehler}}

\newglossaryentry{redundanz}{
  name=Redundanz,
  description={Daten-Anteil, der über die reine Information
    hinausgeht und der Erkennung oder Korrektur von
    Übertragungsfehlern dient. Preis für Robustheit}}

\newglossaryentry{ean13}{
  name=EAN-13,
  description={European Article Number, 13-stelliger Strichcode auf
    fast jedem Konsumgut. Letzte Ziffer ist eine Prüfziffer (Modulo
    10 mit Gewichten 1 und 3)}}

\newglossaryentry{isbn10}{
  name=ISBN-10,
  description={Bis 2007 verwendete 10-stellige Buchnummer. Prüfziffer
    Modulo 11; deshalb gelegentlich „X" als Prüfsymbol nötig}}

\newglossaryentry{luhn}{
  name=Luhn-Algorithmus,
  description={Prüfziffer-Verfahren für Kreditkartennummern (Hans Peter
    Luhn, 1954). Verdoppelt jede zweite Ziffer von rechts und summiert
    die Quersummen}}

\newglossaryentry{hammingdistanz}{
  name=Hamming-Distanz,
  description={Anzahl der Stellen, an denen sich zwei gleich lange
    Wörter unterscheiden. Zentrales Maß für Codes}}

\newglossaryentry{mindestdistanz}{
  name=Mindestdistanz,
  description={Kleinste Hamming-Distanz zwischen je zwei Codewörtern
    eines Codes. Bestimmt die Korrektur-/Erkennungs-Fähigkeit:
    \(d-1\) Fehler erkennbar, \(\lfloor (d-1)/2 \rfloor\)
    korrigierbar}}

\newglossaryentry{code}{
  name=Code,
  description={Menge gültiger Codewörter. Die Hamming-Distanzen
    zwischen ihnen bestimmen die Korrektureigenschaften}}

\newglossaryentry{codewort}{
  name=Codewort,
  description={Element eines Codes, also eine gültig codierte Bit-
    oder Symbolfolge}}

\newglossaryentry{hammingschranke}{
  name=Hamming-Schranke,
  description={Untere Grenze für die Anzahl Bits, die zur
    Einzelfehler-Korrektur nötig sind: \(2^k \cdot (n+1) \le 2^n\).
    Wird vom \((7,4)\)-Hamming-Code mit Gleichheit erreicht („perfekter
    Code")}}

\newglossaryentry{hammingcode}{
  name=Hamming-Code,
  description={Familie linearer Codes (Richard Hamming, 1950), die
    einen Einzel-Bitfehler korrigieren. Der \((7,4)\)-Variante codiert
    4 Datenbits in 7 Codebits}}

\newglossaryentry{secded}{
  name=SECDED,
  description={Single Error Correction, Double Error Detection.
    Erweiterung eines Hamming-Codes um ein Gesamt-Paritätsbit:
    erkennt Doppelfehler zusätzlich zur Einzelfehler-Korrektur. In
    ECC-RAM Standard}}

\newglossaryentry{buendelfehler}{
  name=Bündelfehler,
  description={Mehrere benachbarte Bitfehler, etwa durch Kratzer auf
    einer CD oder Funkfading. Realer als zufällige Einzelfehler;
    motiviert Interleaving}}

\newglossaryentry{interleaving}{
  name=Interleaving,
  description={Verschachtelung von Codewörtern vor der Übertragung,
    sodass ein Bündelfehler aus jedem Codewort nur ein einziges Bit
    trifft. Steckt in jeder CD und vielen Funkstandards}}

\newglossaryentry{crc}{
  name=CRC,
  description={Cyclic Redundancy Check. Standardverfahren der
    Fehlererkennung. Daten werden als Polynom interpretiert und
    modulo eines Generatorpolynoms gerechnet}}

\newglossaryentry{generatorpolynom}{
  name=Generatorpolynom,
  description={Festes Polynom \(G(x)\) bei CRC oder Reed-Solomon, das
    die Code-Eigenschaften bestimmt. Codewörter sind durch \(G(x)\)
    teilbar}}

\newglossaryentry{koerper}{
  name=Körper,
  description={Mathematische Struktur, in der man uneingeschränkt
    addieren, subtrahieren, multiplizieren und dividieren kann
    (Ausnahme: Division durch Null). Beispiele: \(\mathbb{R}\),
    \(\mathbb{Q}\), \(\mathbb{Z}/p\mathbb{Z}\) für \(p\) prim}}

\newglossaryentry{endlicherkoerper}{
  name=Endlicher Körper,
  description={Körper mit endlich vielen Elementen. Tritt nur in
    Größen \(p^n\) auf (\(p\) prim). Notation \(GF(p^n)\) für
    \emph{Galois-Feld}. Für Reed-Solomon im Datamatrix wird
    \(GF(2^8)\) verwendet}}

\newglossaryentry{galoisfeld}{
  name=Galois-Feld,
  description={Synonym für endlicher Körper, benannt nach Évariste
    Galois (1811--1832). Geschrieben \(GF(p^n)\)}}

\newglossaryentry{charakteristik}{
  name=Charakteristik,
  description={Kleinste positive Zahl \(p\), für die \(p \cdot 1 = 0\)
    gilt. In jedem \(GF(2^n)\) ist die Charakteristik 2 — daraus
    folgt der „Freshman's Dream" \((a+b)^2 = a^2 + b^2\)}}

\newglossaryentry{irreduziblespolynom}{
  name=Irreduzibles Polynom,
  description={Polynom, das nicht in zwei Polynome niedrigeren Grades
    zerlegt werden kann. Analoges Konzept zur Primzahl. Wird gebraucht,
    um \(GF(2^n)\) als Körper zu bauen}}

\newglossaryentry{inverses}{
  name=Multiplikatives Inverses,
  description={Zu einem Element \(a \ne 0\) eines Körpers das
    Element \(a^{-1}\) mit \(a \cdot a^{-1} = 1\). Existiert in
    Körpern für jedes Element außer Null}}

\newglossaryentry{stuetzstelle}{
  name=Stützstelle,
  description={Punkt \(x_i\), an dem ein Polynom \(f\) ausgewertet
    wird, um den Wert \(f(x_i)\) zu erhalten. Bei Reed-Solomon werden
    \(n\) Stützstellen verwendet, von denen \(k\) zur Rekonstruktion
    reichen}}

\newglossaryentry{reedsolomon}{
  name=Reed-Solomon-Code,
  description={Fehlerkorrigierender Code über \(GF(2^n)\). Daten
    werden als Polynom interpretiert und an \(n\) Stützstellen
    ausgewertet; die \(n\) Werte sind das Codewort. Korrigiert bis
    zu \(\lfloor (n-k)/2 \rfloor\) Symbolfehler}}

\newglossaryentry{aushloeschung}{
  name=Auslöschung,
  description={Fehler mit bekannter Position (engl.\ \emph{erasure}).
    Im Gegensatz zum allgemeinen Fehler („wir wissen nicht, welches
    Symbol verfälscht ist") doppelt so günstig zu korrigieren}}

\newglossaryentry{horner}{
  name=Horner-Schema,
  description={Effizientes Verfahren zur Polynom-Auswertung:
    \(a_0 + x(a_1 + x(a_2 + \dots))\). Spart Multiplikationen
    gegenüber der naiven Auswertung}}


% Defaults, falls ein Master vergisst sie zu setzen.
\providecommand{\maxetappe}{8}
\providecommand{\hauptbuchtitelseite}{}

\begin{document}

\frontmatter
\hauptbuchtitelseite
\pagestyle{plain}

\chapter*{Vorwort}
\addcontentsline{toc}{chapter}{Vorwort}

\section*{Worum es geht}

Auf jeder Schokoladenpackung, in jedem Funkpaket, auf jeder
zerkratzten CD und in jedem Datamatrix-Code auf einem Versandetikett
arbeitet im Hintergrund \emph{Mathematik}, die Fehler erkennt und
korrigiert. Dieses Buch baut diese Mathematik Stück für Stück auf —
von der einfachsten Prüfziffer bis zum vollständigen Datamatrix-Code
mit Reed-Solomon.

\section*{Über Greta}

Das Buch ist als zweiwöchiges Praktikum gedacht und richtet sich
direkt an „Greta", eine Schülerin der zehnten Klasse. Greta gibt es
wirklich; aber selbst wenn du keine Schülerin namens Greta bist, kannst
du das Buch als Greta lesen. Der direkte Ton ist Absicht: niemand
erklärt etwas Schwieriges in der dritten Person.

Vorwissen, das wir voraussetzen: Schulmathematik bis Klasse 10
(Bruchrechnen, Funktionen, ein bisschen Beweise mitdenken) und etwas
Python-Erfahrung (\texttt{for}-Schleifen, Listen, Funktionen — oder die
Bereitschaft, das nebenbei aufzuholen).

\section*{Wie das Buch aufgebaut ist}

Jedes Kapitel ist ein \emph{Tag}. Ein Tag dauert ungefähr zwei Stunden
und besteht aus mehreren Blöcken mit Zeitangaben. In jedem Block
wechseln sich drei Aufgabentypen ab, die du am Strich und Symbol auf
der Seitenaußenkante erkennst:

\begin{itemize}
\item \textcolor{cBleistift}{\textbf{Bleistift-Übungen}} (orange) — mit
  Stift und Karopapier. Selbstdenken vor Selbstrechnen vor Selbstprüfen.
\item \textcolor{cPython}{\textbf{Python-Einheiten}} (blau) — Code in
  Thonny abtippen, ausführen, durchspielen. Computer brute-forcen
  schneller als wir; das ist eines unserer Lieblingsspielzeuge.
\item \textcolor{cWerkzeug}{\textbf{Werkzeug-Checks}} (grün) — kurze
  Setup-Schritte, damit du nicht beim Installieren von Thonny stecken
  bleibst.
\end{itemize}

Im Anhang stehen die \textbf{Lösungen} (Anhang A) und ein
\textbf{Glossar} aller Fachbegriffe. Lösungen sind absichtlich
ausführlich; den Aha-Moment hat man trotzdem nur, wenn man sie erst
nach eigenem Versuch liest.

\section*{Material}

\begin{itemize}
\item Bleistift, Radiergummi, kariertes Papier
\item drei Buntstifte für die Diagramme an Tag 3
\item ein Laptop mit Python (Thonny ist die empfohlene Umgebung)
\item ein Produkt mit Strichcode (Schokolade, Buch — egal) für Tag 1
\end{itemize}

\section*{Tempo}

Die Tagesabschnitte sind als Häppchen gedacht. Wenn ein Tag sich nach
zwei Stunden noch nicht zu Ende anfühlt, ist es kein Verbrechen, ihn
auf zwei Sitzungen aufzuteilen. Tag 5 (endliche Körper) ist
konzeptionell der schwerste — gerade dort darf Pause sein.

Viel Spaß beim Lesen, Rechnen und Tippen.

\tableofcontents

\mainmatter
\pagestyle{scrheadings}

% Etappen 1 .. \maxetappe der Reihe nach einbinden. \include sorgt
% pro Datei für eigenen aux-File und Pagebreak.
\foreach \etappennr in {1,...,\maxetappe}{%
  \include{kapitel/tag\etappennr}%
}

\appendix

% Kapitel-spezifische Anhänge (z. B. Tabellen, Werkstatt-Material).
% Pro Etappe wird anhang/tagN.tex eingebunden, wenn die Datei existiert.
\foreach \etappennr in {1,...,\maxetappe}{%
  \IfFileExists{anhang/tag\etappennr.tex}{%
    \input{anhang/tag\etappennr}%
  }{}%
}

% Lösungen zu allen bisher behandelten Etappen, gesammelt im Anhang.
\chapter{Lösungen}

\begin{quote}
Erst nach eigenem Versuch ansehen — der Aha-Moment beim Selber-Lösen
ist wichtiger als die fertige Lösung am Ende.
\end{quote}

\foreach \etappennr in {1,...,\maxetappe}{%
  \input{loesungen/tag\etappennr}%
}

% \glsaddall zeigt alle Glossar-Einträge (auch die, die im Text nicht
% per \gls{...} referenziert wurden). Solange wir die Markups nicht
% systematisch durch alle Kapitel gezogen haben, ist das die
% praktikable Variante.
\glsaddall
\printglossary[title=Glossar, toctitle=Glossar]

\printindex

\end{document}
 auf den hier definierten Body zu.
% =============================================================

% =============================================================
% Präambel — Pakete, Schriften, Geometrie, Sprache, Kopf-/Fußzeile.
%
% Wird von buch.tex eingebunden, NACH \documentclass.
% Layout-Bausteine (Aufgaben-Umgebungen) sind in befehle.tex,
% Farben in farben.tex.
% =============================================================

% --- Geometrie nach doku/LAYOUT_BRIEFING.md ---
\usepackage[a4paper,
            inner=2.5cm, outer=4.5cm,
            top=2.0cm, bottom=2.5cm,
            marginparwidth=2.7cm, marginparsep=0.15cm,
            headheight=15pt]{geometry}

% --- Schriften ---
% Source Pro Familie als Hausschriften, vendored unter latex/fonts/
% (siehe latex/fonts/README.md für Provenance). Source Serif für den
% Fließtext, Source Sans für Überschriften/Marginalien, Source Code
% für Quelltexte. Math-Formeln werden weiter unten via unicode-math
% gesetzt (STIX Two Math passt visuell zu Source Serif).
\usepackage{fontspec}
\setmainfont{SourceSerif4}[
  Path=fonts/source-serif/,
  Extension=.otf,
  UprightFont=*-Regular,
  ItalicFont=*-It,
  BoldFont=*-Bold,
  BoldItalicFont=*-BoldIt,
  Scale=1.00,
]
\setsansfont{SourceSans3}[
  Path=fonts/source-sans/,
  Extension=.otf,
  UprightFont=*-Regular,
  ItalicFont=*-It,
  BoldFont=*-Bold,
  BoldItalicFont=*-BoldIt,
  Scale=0.95,
]
\setmonofont{SourceCodePro}[
  Path=fonts/source-code-pro/,
  Extension=.otf,
  UprightFont=*-Regular,
  ItalicFont=*-It,
  BoldFont=*-Bold,
  BoldItalicFont=*-BoldIt,
  Scale=0.88,
]
% Marginalien-Symbole (✏ ⌨ ⚙) brauchen einen Symbol-fähigen Font.
% Source Sans hat die Glyphen nicht; wir behalten DejaVu Sans als
% Fallback nur dafür (System-Font, im Container via fonts-dejavu).
\newfontfamily\markerfont{DejaVu Sans}

% --- Sprache, Mikrotypografie ---
\usepackage[ngerman]{babel}
% babel-ngerman macht ASCII-" zum Shorthand-Aktivator ("u → ü, "f → ff
% etc.). Das hat schon zu absurden Verschmelzungen wie "echterFFehler"
% statt "echter" Fehler" geführt. Wir schalten den Shorthand ab; für
% typografisch korrekte deutsche Anführungszeichen nutzen wir Unicode
% („…") direkt im Quelltext.
\shorthandoff{"}
\usepackage{microtype}

% --- Mathematik (\bmod, \le, \ge, \checkmark, …) ---
% amsmath für Mathematik-Umgebungen (align, gather, …).
% unicode-math + STIX Two Math ersetzt die Computer-Modern-Math-
% Schriften; Stilfamilie passt visuell zu Source Serif.
\usepackage{amsmath}
\usepackage{unicode-math}
\setmathfont{STIXTwoMath-Regular.otf}[
  Path=fonts/stix-two/,
]

% --- Glossar (Polish-Sprint M6) ---
% automake=immediate ruft makeglossaries per \write18 auf (shell-escape
% ist eh schon aktiv wegen minted), nonumberlist unterdrückt die
% Seitenzahlen hinter den Glossar-Einträgen (sonst stehen dort die
% Seiten, an denen \gls genutzt wurde).
\usepackage[automake=immediate, toc, nonumberlist]{glossaries-extra}
\makeglossaries

% --- Sachindex ---
\usepackage{imakeidx}
\makeindex[intoc, columns=2]

% --- Farben & Hyperlinks ---
\usepackage{xcolor}
\usepackage[hidelinks]{hyperref}
% pdfpagelabels ist hyperref-Default; wir setzen nur plainpages=false,
% damit hyperref beim \pagenumbering-Wechsel Frontmatter→Mainmatter
% keine Duplikat-Warnung wirft und PDF-Reader die Buchseiten-Labels
% (i, ii, …, 1, 2, …) korrekt anzeigen.
\hypersetup{plainpages=false}

% --- Layout-Bausteine ---
% tcolorbox: nur die Libraries laden, die wir tatsächlich brauchen
% (skins für `enhanced` und `borderline`, breakable für mehrseitige Boxen).
% [most] würde tikzfill u.a. mitziehen, was wir nicht nutzen.
\usepackage{tcolorbox}
\tcbuselibrary{skins, breakable}
\usepackage{changepage}
% \strictpagecheck macht \checkoddpage zuverlässig auch bei Page-Breaks
% (sonst kann die Seitenparität pro Page-Part falsch geprüft werden).
\strictpagecheck

% --- TikZ (für Venn-Diagramme etc.) ---
\usepackage{tikz}
\usetikzlibrary{positioning, calc}

% --- Tabellen ---
\usepackage{booktabs, array, longtable}

% --- Listen mit (a)/(b)/(c)-Labels für Aufgabenteile ---
\usepackage{enumitem}

% --- needspace: erzwingt Mindestplatz vor Aufgabenblöcken ---
\usepackage{needspace}

% --- Code-Highlighting via minted ---
% minted ruft Pygments per Shell auf; deshalb -shell-escape in .latexmkrc.
\usepackage{minted}
% Hintergrundfarbe für Code-Blöcke (sehr helles Grau, naive CMYK aus
% RGB(246, 247, 248), entspricht der bisherigen pandoc-shadecolor).
\definecolor{codebg}{cmyk}{0.01, 0.00, 0.00, 0.03}
\setminted{
  style=tango,
  fontsize=\small,
  bgcolor=codebg,
  frame=none,
  breaklines=true,
  autogobble=true,
  linenos=false,
  baselinestretch=1.05,
  xleftmargin=2pt,
  xrightmargin=2pt,
}

% --- Kopf-/Fußzeile via scrlayer-scrpage ---
\usepackage[autooneside=false, automark]{scrlayer-scrpage}
\clearpairofpagestyles
\automark[chapter]{chapter}
\ihead{\small\textit{\headmark}}
\ohead{\small\pagemark}
\KOMAoptions{headsepline=0.3pt}

% --- Witwen / Waisen vermeiden ---
\widowpenalty=10000
\clubpenalty=10000

% --- Etwas Toleranz für die schmaleren Spalten ---
\setlength{\emergencystretch}{3em}

% =============================================================
% Farbset 3 aus doku/LAYOUT_BRIEFING.md (Pastell, hell).
%
% Definiert in CMYK, weil das Buch in den Druck soll.
% Die ursprünglichen RGB-Werte aus dem Briefing stehen als Kommentar
% darüber; die CMYK-Werte sind naiv per Standardformel umgerechnet
% (K = 1 - max(R,G,B)/255 etc.). Für endgültige Druckqualität ggf.
% mit ICC-Profil (z. B. ISOcoated_v2_eci) nachjustieren.
%
% Diese drei Farben werden konsequent verwendet:
%   - für den Strich am Aufgabenblock
%   - für das Symbol in der Marginalie
%   - für die Kapitälchen-Bezeichnung unter dem Symbol
% =============================================================

% Bleistift — RGB(235, 160,  80), helles Orange
\definecolor{cBleistift}{cmyk}{0.00, 0.32, 0.66, 0.08}

% Python — RGB( 95, 140, 180), sanftes Blau
\definecolor{cPython}   {cmyk}{0.47, 0.22, 0.00, 0.29}

% Werkzeug — RGB(185, 200, 110), helles Lindgrün
\definecolor{cWerkzeug} {cmyk}{0.08, 0.00, 0.45, 0.22}

% =============================================================
% Aufgaben-Umgebungen.
%
% Drei Komfort-Umgebungen für die drei Aufgabentypen:
%   \begin{bleistiftuebung}{N}{Titel}  ...  \end{bleistiftuebung}
%   \begin{pythoneinheit} {N}{Titel}  ...  \end{pythoneinheit}
%   \begin{werkzeugcheck}    {Titel}  ...  \end{werkzeugcheck}
%
% Jede Umgebung erzeugt:
%   - eine farbige tcolorbox um den Aufgabentext, mit einem 2,5-pt-Strich
%     am AUSSEN-Rand der Seite (recto = rechts, verso = links),
%   - eine Marginalie auf derselben Außenseite mit Symbol und
%     Kapitälchen-Bezeichnung in der Aufgaben-Farbe.
%
% Die Bausteine stammen aus doku/LAYOUT_BRIEFING.md, dort steht auch
% die Begründung für jede Designentscheidung.
% =============================================================

% --- aufgabebox: tcolorbox mit page-aware Strich am Außenrand ---
% Eine einzige Box, kein Recto/Verso-Split mehr. Der overlay-Hook wird
% bei jedem Page-Part neu ausgewertet (auch wenn die Aufgabe per
% breakable über Seiten geht), prüft die aktuelle Seitenparität und
% zeichnet den Strich auf der jeweiligen Außenseite.
%
% left/right=8pt sorgt für Innenabstand zum Strich auf beiden Seiten;
% der Text bleibt immer mittig in der Box.
\newtcolorbox{aufgabebox}[1]{%
  enhanced, breakable,
  colback=white, colframe=white,
  boxrule=0pt, arc=0pt, outer arc=0pt,
  boxsep=0pt,
  left=8pt, right=8pt, top=4pt, bottom=4pt,
  before skip=0.6em, after skip=0.6em,
  overlay={%
    \checkoddpage
    \ifoddpage
      \draw[line width=2.5pt, color=#1]
        (frame.north east) -- (frame.south east);%
    \else
      \draw[line width=2.5pt, color=#1]
        (frame.north west) -- (frame.south west);%
    \fi
  },
}

% --- Marginalien-Inhalt für Recto und Verso ---
% Auf Recto-Seiten (Marginalspalte rechts): linksbündig zum Strich.
% Auf Verso-Seiten (Marginalspalte links):  rechtsbündig zum Strich.
% Das \vspace*{12pt} kompensiert das `before skip` und das `top` der
% tcolorbox, damit das Symbol mit der ersten Zeile der Aufgabenbox fluchtet.
\newcommand{\markercontentRecto}[4]{%
  \vspace*{12pt}\raggedright\color{#4}%
  {\markerfont\fontsize{24}{26}\selectfont #1}\\[0.05em]
  {\small\textsc{#2}}\\
  {\small\textsc{#3}}%
}
\newcommand{\markercontentVerso}[4]{%
  \vspace*{12pt}\raggedleft\color{#4}%
  {\markerfont\fontsize{24}{26}\selectfont #1}\\[0.05em]
  {\small\textsc{#2}}\\
  {\small\textsc{#3}}%
}

% --- Aufgaben-Umgebung ---
% Marginpar wird mit beiden Inhaltsvarianten aufgerufen (recto/verso);
% LaTeX wählt selbst die passende, je nachdem auf welcher Seite die
% Aufgabe beginnt. Der page-aware Strich entsteht im Overlay der
% aufgabebox und wechselt automatisch beim Seitenumbruch mit.
\newenvironment{aufgabe}[4]%
  {% #1 = Symbol, #2 = Label-Zeile-1, #3 = Label-Zeile-2, #4 = Farbe
   \marginpar[\markercontentVerso{#1}{#2}{#3}{#4}]%
              {\markercontentRecto{#1}{#2}{#3}{#4}}%
   \begin{aufgabebox}{#4}\ignorespaces}
  {\end{aufgabebox}}

% --- Komfort-Wrapper für die drei Aufgabentypen ---
% Konvention: der erste Teil der Aufgabe ist der Titel (fett), gefolgt
% vom Aufgabentext. Der Titel wird hier als Argument übergeben.
%
% Wir nutzen {\bfseries ...\par} statt \textbf{...\par}, weil \textbf
% als Inline-Befehl keinen Absatzumbruch (\par) im Argument verträgt
% ("Paragraph ended before \text@command was complete").
\newenvironment{bleistiftuebung}[2]%
  {\begin{aufgabe}{✏}{Bleistift-}{Übung #1}{cBleistift}{\bfseries #2.\par}}
  {\end{aufgabe}}

\newenvironment{pythoneinheit}[2]%
  {\begin{aufgabe}{⌨}{Python-}{Einheit #1}{cPython}{\bfseries #2.\par}}
  {\end{aufgabe}}

\newenvironment{werkzeugcheck}[1]%
  {\begin{aufgabe}{⚙}{Werkzeug-}{Check}{cWerkzeug}{\bfseries #1.\par}}
  {\end{aufgabe}}

% --- Code-Snippet-Einbindung ---
% \pythoncode{tagN/file.py} bindet eine Python-Datei vollständig ein.
% Pfad ist relativ zu latex/code/. Die .py-Dateien sind eigenständig
% lauffähig, sodass Greta sie 1:1 in Thonny verwenden kann.
\newcommand{\pythoncode}[1]{\inputminted{python}{code/#1}}

% \pythoncodeteil{tagN/file.py}{regionname} zeigt nur eine Region
% aus einer Datei. Region-Marker in der Quelle:
%     # region <name>
%     ... Code ...
%     # endregion
% Greta tippt die Datei sukzessive auf, im Buch zeigen wir pro
% Python-Einheit jeweils nur das aktuell relevante Stück.
%
% Implementierung: extract_region.py schneidet die Region zur
% Build-Zeit in eine temporäre Datei, die minted dann einbindet.
% Der Counter sorgt dafür, dass jede Einbindung ihre eigene Datei
% bekommt — nötig, weil minted das Snippet zur Build-Zeit cached.
\newcounter{snippetcntr}
\makeatletter
\newcommand{\pythoncodeteil}[2]{%
  \stepcounter{snippetcntr}%
  \protected@edef\@snippetfile{.snippets/snippet-\arabic{snippetcntr}.py}%
  \immediate\write18{python3 scripts/extract_region.py code/#1 #2 > \@snippetfile}%
  \inputminted{python}{\@snippetfile}%
}
\makeatother

% =============================================================
% Glossar-Einträge.
%
% Eingebunden in der Präambel (vor \begin{document}). Im Text
% verwendet via \gls{schluessel} an der ersten Erwähnung; weitere
% Auftritte können auch \gls{} sein, der Glossar-Eintrag wird dann
% nur noch normal gesetzt.
%
% Schlüssel sind kurz und sprechend; Anzeigename steht in `name=`.
% =============================================================

\newglossaryentry{modulo}{
  name=Modulo,
  description={Operation \(a \bmod n\): der Rest, der bei der Division
    von \(a\) durch \(n\) bleibt. Universalwerkzeug der
    Codierungstheorie}}

\newglossaryentry{paritaet}{
  name=Parität,
  description={Eigenschaft einer Bitfolge, die sich aus der Anzahl der
    Einsen ergibt: \emph{gerade Parität}, wenn die Anzahl gerade
    ist, sonst \emph{ungerade}. Eine Paritätsprüfung erkennt jede
    ungerade Anzahl gekippter Bits}}

\newglossaryentry{paritaetsbit}{
  name=Paritätsbit,
  description={Zusätzliches Bit, das so gewählt wird, dass die
    Gesamt-Parität (meist gerade) zustande kommt. Erkennt
    Einzel-Bitfehler}}

\newglossaryentry{redundanz}{
  name=Redundanz,
  description={Daten-Anteil, der über die reine Information
    hinausgeht und der Erkennung oder Korrektur von
    Übertragungsfehlern dient. Preis für Robustheit}}

\newglossaryentry{ean13}{
  name=EAN-13,
  description={European Article Number, 13-stelliger Strichcode auf
    fast jedem Konsumgut. Letzte Ziffer ist eine Prüfziffer (Modulo
    10 mit Gewichten 1 und 3)}}

\newglossaryentry{isbn10}{
  name=ISBN-10,
  description={Bis 2007 verwendete 10-stellige Buchnummer. Prüfziffer
    Modulo 11; deshalb gelegentlich „X" als Prüfsymbol nötig}}

\newglossaryentry{luhn}{
  name=Luhn-Algorithmus,
  description={Prüfziffer-Verfahren für Kreditkartennummern (Hans Peter
    Luhn, 1954). Verdoppelt jede zweite Ziffer von rechts und summiert
    die Quersummen}}

\newglossaryentry{hammingdistanz}{
  name=Hamming-Distanz,
  description={Anzahl der Stellen, an denen sich zwei gleich lange
    Wörter unterscheiden. Zentrales Maß für Codes}}

\newglossaryentry{mindestdistanz}{
  name=Mindestdistanz,
  description={Kleinste Hamming-Distanz zwischen je zwei Codewörtern
    eines Codes. Bestimmt die Korrektur-/Erkennungs-Fähigkeit:
    \(d-1\) Fehler erkennbar, \(\lfloor (d-1)/2 \rfloor\)
    korrigierbar}}

\newglossaryentry{code}{
  name=Code,
  description={Menge gültiger Codewörter. Die Hamming-Distanzen
    zwischen ihnen bestimmen die Korrektureigenschaften}}

\newglossaryentry{codewort}{
  name=Codewort,
  description={Element eines Codes, also eine gültig codierte Bit-
    oder Symbolfolge}}

\newglossaryentry{hammingschranke}{
  name=Hamming-Schranke,
  description={Untere Grenze für die Anzahl Bits, die zur
    Einzelfehler-Korrektur nötig sind: \(2^k \cdot (n+1) \le 2^n\).
    Wird vom \((7,4)\)-Hamming-Code mit Gleichheit erreicht („perfekter
    Code")}}

\newglossaryentry{hammingcode}{
  name=Hamming-Code,
  description={Familie linearer Codes (Richard Hamming, 1950), die
    einen Einzel-Bitfehler korrigieren. Der \((7,4)\)-Variante codiert
    4 Datenbits in 7 Codebits}}

\newglossaryentry{secded}{
  name=SECDED,
  description={Single Error Correction, Double Error Detection.
    Erweiterung eines Hamming-Codes um ein Gesamt-Paritätsbit:
    erkennt Doppelfehler zusätzlich zur Einzelfehler-Korrektur. In
    ECC-RAM Standard}}

\newglossaryentry{buendelfehler}{
  name=Bündelfehler,
  description={Mehrere benachbarte Bitfehler, etwa durch Kratzer auf
    einer CD oder Funkfading. Realer als zufällige Einzelfehler;
    motiviert Interleaving}}

\newglossaryentry{interleaving}{
  name=Interleaving,
  description={Verschachtelung von Codewörtern vor der Übertragung,
    sodass ein Bündelfehler aus jedem Codewort nur ein einziges Bit
    trifft. Steckt in jeder CD und vielen Funkstandards}}

\newglossaryentry{crc}{
  name=CRC,
  description={Cyclic Redundancy Check. Standardverfahren der
    Fehlererkennung. Daten werden als Polynom interpretiert und
    modulo eines Generatorpolynoms gerechnet}}

\newglossaryentry{generatorpolynom}{
  name=Generatorpolynom,
  description={Festes Polynom \(G(x)\) bei CRC oder Reed-Solomon, das
    die Code-Eigenschaften bestimmt. Codewörter sind durch \(G(x)\)
    teilbar}}

\newglossaryentry{koerper}{
  name=Körper,
  description={Mathematische Struktur, in der man uneingeschränkt
    addieren, subtrahieren, multiplizieren und dividieren kann
    (Ausnahme: Division durch Null). Beispiele: \(\mathbb{R}\),
    \(\mathbb{Q}\), \(\mathbb{Z}/p\mathbb{Z}\) für \(p\) prim}}

\newglossaryentry{endlicherkoerper}{
  name=Endlicher Körper,
  description={Körper mit endlich vielen Elementen. Tritt nur in
    Größen \(p^n\) auf (\(p\) prim). Notation \(GF(p^n)\) für
    \emph{Galois-Feld}. Für Reed-Solomon im Datamatrix wird
    \(GF(2^8)\) verwendet}}

\newglossaryentry{galoisfeld}{
  name=Galois-Feld,
  description={Synonym für endlicher Körper, benannt nach Évariste
    Galois (1811--1832). Geschrieben \(GF(p^n)\)}}

\newglossaryentry{charakteristik}{
  name=Charakteristik,
  description={Kleinste positive Zahl \(p\), für die \(p \cdot 1 = 0\)
    gilt. In jedem \(GF(2^n)\) ist die Charakteristik 2 — daraus
    folgt der „Freshman's Dream" \((a+b)^2 = a^2 + b^2\)}}

\newglossaryentry{irreduziblespolynom}{
  name=Irreduzibles Polynom,
  description={Polynom, das nicht in zwei Polynome niedrigeren Grades
    zerlegt werden kann. Analoges Konzept zur Primzahl. Wird gebraucht,
    um \(GF(2^n)\) als Körper zu bauen}}

\newglossaryentry{inverses}{
  name=Multiplikatives Inverses,
  description={Zu einem Element \(a \ne 0\) eines Körpers das
    Element \(a^{-1}\) mit \(a \cdot a^{-1} = 1\). Existiert in
    Körpern für jedes Element außer Null}}

\newglossaryentry{stuetzstelle}{
  name=Stützstelle,
  description={Punkt \(x_i\), an dem ein Polynom \(f\) ausgewertet
    wird, um den Wert \(f(x_i)\) zu erhalten. Bei Reed-Solomon werden
    \(n\) Stützstellen verwendet, von denen \(k\) zur Rekonstruktion
    reichen}}

\newglossaryentry{reedsolomon}{
  name=Reed-Solomon-Code,
  description={Fehlerkorrigierender Code über \(GF(2^n)\). Daten
    werden als Polynom interpretiert und an \(n\) Stützstellen
    ausgewertet; die \(n\) Werte sind das Codewort. Korrigiert bis
    zu \(\lfloor (n-k)/2 \rfloor\) Symbolfehler}}

\newglossaryentry{aushloeschung}{
  name=Auslöschung,
  description={Fehler mit bekannter Position (engl.\ \emph{erasure}).
    Im Gegensatz zum allgemeinen Fehler („wir wissen nicht, welches
    Symbol verfälscht ist") doppelt so günstig zu korrigieren}}

\newglossaryentry{horner}{
  name=Horner-Schema,
  description={Effizientes Verfahren zur Polynom-Auswertung:
    \(a_0 + x(a_1 + x(a_2 + \dots))\). Spart Multiplikationen
    gegenüber der naiven Auswertung}}


% Defaults, falls ein Master vergisst sie zu setzen.
\providecommand{\maxetappe}{8}
\providecommand{\hauptbuchtitelseite}{}

\begin{document}

\frontmatter
\hauptbuchtitelseite
\pagestyle{plain}

\chapter*{Vorwort}
\addcontentsline{toc}{chapter}{Vorwort}

\section*{Worum es geht}

Auf jeder Schokoladenpackung, in jedem Funkpaket, auf jeder
zerkratzten CD und in jedem Datamatrix-Code auf einem Versandetikett
arbeitet im Hintergrund \emph{Mathematik}, die Fehler erkennt und
korrigiert. Dieses Buch baut diese Mathematik Stück für Stück auf —
von der einfachsten Prüfziffer bis zum vollständigen Datamatrix-Code
mit Reed-Solomon.

\section*{Über Greta}

Das Buch ist als zweiwöchiges Praktikum gedacht und richtet sich
direkt an „Greta", eine Schülerin der zehnten Klasse. Greta gibt es
wirklich; aber selbst wenn du keine Schülerin namens Greta bist, kannst
du das Buch als Greta lesen. Der direkte Ton ist Absicht: niemand
erklärt etwas Schwieriges in der dritten Person.

Vorwissen, das wir voraussetzen: Schulmathematik bis Klasse 10
(Bruchrechnen, Funktionen, ein bisschen Beweise mitdenken) und etwas
Python-Erfahrung (\texttt{for}-Schleifen, Listen, Funktionen — oder die
Bereitschaft, das nebenbei aufzuholen).

\section*{Wie das Buch aufgebaut ist}

Jedes Kapitel ist ein \emph{Tag}. Ein Tag dauert ungefähr zwei Stunden
und besteht aus mehreren Blöcken mit Zeitangaben. In jedem Block
wechseln sich drei Aufgabentypen ab, die du am Strich und Symbol auf
der Seitenaußenkante erkennst:

\begin{itemize}
\item \textcolor{cBleistift}{\textbf{Bleistift-Übungen}} (orange) — mit
  Stift und Karopapier. Selbstdenken vor Selbstrechnen vor Selbstprüfen.
\item \textcolor{cPython}{\textbf{Python-Einheiten}} (blau) — Code in
  Thonny abtippen, ausführen, durchspielen. Computer brute-forcen
  schneller als wir; das ist eines unserer Lieblingsspielzeuge.
\item \textcolor{cWerkzeug}{\textbf{Werkzeug-Checks}} (grün) — kurze
  Setup-Schritte, damit du nicht beim Installieren von Thonny stecken
  bleibst.
\end{itemize}

Im Anhang stehen die \textbf{Lösungen} (Anhang A) und ein
\textbf{Glossar} aller Fachbegriffe. Lösungen sind absichtlich
ausführlich; den Aha-Moment hat man trotzdem nur, wenn man sie erst
nach eigenem Versuch liest.

\section*{Material}

\begin{itemize}
\item Bleistift, Radiergummi, kariertes Papier
\item drei Buntstifte für die Diagramme an Tag 3
\item ein Laptop mit Python (Thonny ist die empfohlene Umgebung)
\item ein Produkt mit Strichcode (Schokolade, Buch — egal) für Tag 1
\end{itemize}

\section*{Tempo}

Die Tagesabschnitte sind als Häppchen gedacht. Wenn ein Tag sich nach
zwei Stunden noch nicht zu Ende anfühlt, ist es kein Verbrechen, ihn
auf zwei Sitzungen aufzuteilen. Tag 5 (endliche Körper) ist
konzeptionell der schwerste — gerade dort darf Pause sein.

Viel Spaß beim Lesen, Rechnen und Tippen.

\tableofcontents

\mainmatter
\pagestyle{scrheadings}

% Etappen 1 .. \maxetappe der Reihe nach einbinden. \include sorgt
% pro Datei für eigenen aux-File und Pagebreak.
\foreach \etappennr in {1,...,\maxetappe}{%
  \include{kapitel/tag\etappennr}%
}

\appendix

% Kapitel-spezifische Anhänge (z. B. Tabellen, Werkstatt-Material).
% Pro Etappe wird anhang/tagN.tex eingebunden, wenn die Datei existiert.
\foreach \etappennr in {1,...,\maxetappe}{%
  \IfFileExists{anhang/tag\etappennr.tex}{%
    \input{anhang/tag\etappennr}%
  }{}%
}

% Lösungen zu allen bisher behandelten Etappen, gesammelt im Anhang.
\chapter{Lösungen}

\begin{quote}
Erst nach eigenem Versuch ansehen — der Aha-Moment beim Selber-Lösen
ist wichtiger als die fertige Lösung am Ende.
\end{quote}

\foreach \etappennr in {1,...,\maxetappe}{%
  \input{loesungen/tag\etappennr}%
}

% \glsaddall zeigt alle Glossar-Einträge (auch die, die im Text nicht
% per \gls{...} referenziert wurden). Solange wir die Markups nicht
% systematisch durch alle Kapitel gezogen haben, ist das die
% praktikable Variante.
\glsaddall
\printglossary[title=Glossar, toctitle=Glossar]

\printindex

\end{document}
 auf den hier definierten Body zu.
% =============================================================

% =============================================================
% Präambel — Pakete, Schriften, Geometrie, Sprache, Kopf-/Fußzeile.
%
% Wird von buch.tex eingebunden, NACH \documentclass.
% Layout-Bausteine (Aufgaben-Umgebungen) sind in befehle.tex,
% Farben in farben.tex.
% =============================================================

% --- Geometrie nach doku/LAYOUT_BRIEFING.md ---
\usepackage[a4paper,
            inner=2.5cm, outer=4.5cm,
            top=2.0cm, bottom=2.5cm,
            marginparwidth=2.7cm, marginparsep=0.15cm,
            headheight=15pt]{geometry}

% --- Schriften ---
% Source Pro Familie als Hausschriften, vendored unter latex/fonts/
% (siehe latex/fonts/README.md für Provenance). Source Serif für den
% Fließtext, Source Sans für Überschriften/Marginalien, Source Code
% für Quelltexte. Math-Formeln werden weiter unten via unicode-math
% gesetzt (STIX Two Math passt visuell zu Source Serif).
\usepackage{fontspec}
\setmainfont{SourceSerif4}[
  Path=fonts/source-serif/,
  Extension=.otf,
  UprightFont=*-Regular,
  ItalicFont=*-It,
  BoldFont=*-Bold,
  BoldItalicFont=*-BoldIt,
  Scale=1.00,
]
\setsansfont{SourceSans3}[
  Path=fonts/source-sans/,
  Extension=.otf,
  UprightFont=*-Regular,
  ItalicFont=*-It,
  BoldFont=*-Bold,
  BoldItalicFont=*-BoldIt,
  Scale=0.95,
]
\setmonofont{SourceCodePro}[
  Path=fonts/source-code-pro/,
  Extension=.otf,
  UprightFont=*-Regular,
  ItalicFont=*-It,
  BoldFont=*-Bold,
  BoldItalicFont=*-BoldIt,
  Scale=0.88,
]
% Marginalien-Symbole (✏ ⌨ ⚙) brauchen einen Symbol-fähigen Font.
% Source Sans hat die Glyphen nicht; wir behalten DejaVu Sans als
% Fallback nur dafür (System-Font, im Container via fonts-dejavu).
\newfontfamily\markerfont{DejaVu Sans}

% --- Sprache, Mikrotypografie ---
\usepackage[ngerman]{babel}
% babel-ngerman macht ASCII-" zum Shorthand-Aktivator ("u → ü, "f → ff
% etc.). Das hat schon zu absurden Verschmelzungen wie "echterFFehler"
% statt "echter" Fehler" geführt. Wir schalten den Shorthand ab; für
% typografisch korrekte deutsche Anführungszeichen nutzen wir Unicode
% („…") direkt im Quelltext.
\shorthandoff{"}
\usepackage{microtype}

% --- Mathematik (\bmod, \le, \ge, \checkmark, …) ---
% amsmath für Mathematik-Umgebungen (align, gather, …).
% unicode-math + STIX Two Math ersetzt die Computer-Modern-Math-
% Schriften; Stilfamilie passt visuell zu Source Serif.
\usepackage{amsmath}
\usepackage{unicode-math}
\setmathfont{STIXTwoMath-Regular.otf}[
  Path=fonts/stix-two/,
]

% --- Glossar (Polish-Sprint M6) ---
% automake=immediate ruft makeglossaries per \write18 auf (shell-escape
% ist eh schon aktiv wegen minted), nonumberlist unterdrückt die
% Seitenzahlen hinter den Glossar-Einträgen (sonst stehen dort die
% Seiten, an denen \gls genutzt wurde).
\usepackage[automake=immediate, toc, nonumberlist]{glossaries-extra}
\makeglossaries

% --- Sachindex ---
\usepackage{imakeidx}
\makeindex[intoc, columns=2]

% --- Farben & Hyperlinks ---
\usepackage{xcolor}
\usepackage[hidelinks]{hyperref}
% pdfpagelabels ist hyperref-Default; wir setzen nur plainpages=false,
% damit hyperref beim \pagenumbering-Wechsel Frontmatter→Mainmatter
% keine Duplikat-Warnung wirft und PDF-Reader die Buchseiten-Labels
% (i, ii, …, 1, 2, …) korrekt anzeigen.
\hypersetup{plainpages=false}

% --- Layout-Bausteine ---
% tcolorbox: nur die Libraries laden, die wir tatsächlich brauchen
% (skins für `enhanced` und `borderline`, breakable für mehrseitige Boxen).
% [most] würde tikzfill u.a. mitziehen, was wir nicht nutzen.
\usepackage{tcolorbox}
\tcbuselibrary{skins, breakable}
\usepackage{changepage}
% \strictpagecheck macht \checkoddpage zuverlässig auch bei Page-Breaks
% (sonst kann die Seitenparität pro Page-Part falsch geprüft werden).
\strictpagecheck

% --- TikZ (für Venn-Diagramme etc.) ---
\usepackage{tikz}
\usetikzlibrary{positioning, calc}

% --- Tabellen ---
\usepackage{booktabs, array, longtable}

% --- Listen mit (a)/(b)/(c)-Labels für Aufgabenteile ---
\usepackage{enumitem}

% --- needspace: erzwingt Mindestplatz vor Aufgabenblöcken ---
\usepackage{needspace}

% --- Code-Highlighting via minted ---
% minted ruft Pygments per Shell auf; deshalb -shell-escape in .latexmkrc.
\usepackage{minted}
% Hintergrundfarbe für Code-Blöcke (sehr helles Grau, naive CMYK aus
% RGB(246, 247, 248), entspricht der bisherigen pandoc-shadecolor).
\definecolor{codebg}{cmyk}{0.01, 0.00, 0.00, 0.03}
\setminted{
  style=tango,
  fontsize=\small,
  bgcolor=codebg,
  frame=none,
  breaklines=true,
  autogobble=true,
  linenos=false,
  baselinestretch=1.05,
  xleftmargin=2pt,
  xrightmargin=2pt,
}

% --- Kopf-/Fußzeile via scrlayer-scrpage ---
\usepackage[autooneside=false, automark]{scrlayer-scrpage}
\clearpairofpagestyles
\automark[chapter]{chapter}
\ihead{\small\textit{\headmark}}
\ohead{\small\pagemark}
\KOMAoptions{headsepline=0.3pt}

% --- Witwen / Waisen vermeiden ---
\widowpenalty=10000
\clubpenalty=10000

% --- Etwas Toleranz für die schmaleren Spalten ---
\setlength{\emergencystretch}{3em}

% =============================================================
% Farbset 3 aus doku/LAYOUT_BRIEFING.md (Pastell, hell).
%
% Definiert in CMYK, weil das Buch in den Druck soll.
% Die ursprünglichen RGB-Werte aus dem Briefing stehen als Kommentar
% darüber; die CMYK-Werte sind naiv per Standardformel umgerechnet
% (K = 1 - max(R,G,B)/255 etc.). Für endgültige Druckqualität ggf.
% mit ICC-Profil (z. B. ISOcoated_v2_eci) nachjustieren.
%
% Diese drei Farben werden konsequent verwendet:
%   - für den Strich am Aufgabenblock
%   - für das Symbol in der Marginalie
%   - für die Kapitälchen-Bezeichnung unter dem Symbol
% =============================================================

% Bleistift — RGB(235, 160,  80), helles Orange
\definecolor{cBleistift}{cmyk}{0.00, 0.32, 0.66, 0.08}

% Python — RGB( 95, 140, 180), sanftes Blau
\definecolor{cPython}   {cmyk}{0.47, 0.22, 0.00, 0.29}

% Werkzeug — RGB(185, 200, 110), helles Lindgrün
\definecolor{cWerkzeug} {cmyk}{0.08, 0.00, 0.45, 0.22}

% =============================================================
% Aufgaben-Umgebungen.
%
% Drei Komfort-Umgebungen für die drei Aufgabentypen:
%   \begin{bleistiftuebung}{N}{Titel}  ...  \end{bleistiftuebung}
%   \begin{pythoneinheit} {N}{Titel}  ...  \end{pythoneinheit}
%   \begin{werkzeugcheck}    {Titel}  ...  \end{werkzeugcheck}
%
% Jede Umgebung erzeugt:
%   - eine farbige tcolorbox um den Aufgabentext, mit einem 2,5-pt-Strich
%     am AUSSEN-Rand der Seite (recto = rechts, verso = links),
%   - eine Marginalie auf derselben Außenseite mit Symbol und
%     Kapitälchen-Bezeichnung in der Aufgaben-Farbe.
%
% Die Bausteine stammen aus doku/LAYOUT_BRIEFING.md, dort steht auch
% die Begründung für jede Designentscheidung.
% =============================================================

% --- aufgabebox: tcolorbox mit page-aware Strich am Außenrand ---
% Eine einzige Box, kein Recto/Verso-Split mehr. Der overlay-Hook wird
% bei jedem Page-Part neu ausgewertet (auch wenn die Aufgabe per
% breakable über Seiten geht), prüft die aktuelle Seitenparität und
% zeichnet den Strich auf der jeweiligen Außenseite.
%
% left/right=8pt sorgt für Innenabstand zum Strich auf beiden Seiten;
% der Text bleibt immer mittig in der Box.
\newtcolorbox{aufgabebox}[1]{%
  enhanced, breakable,
  colback=white, colframe=white,
  boxrule=0pt, arc=0pt, outer arc=0pt,
  boxsep=0pt,
  left=8pt, right=8pt, top=4pt, bottom=4pt,
  before skip=0.6em, after skip=0.6em,
  overlay={%
    \checkoddpage
    \ifoddpage
      \draw[line width=2.5pt, color=#1]
        (frame.north east) -- (frame.south east);%
    \else
      \draw[line width=2.5pt, color=#1]
        (frame.north west) -- (frame.south west);%
    \fi
  },
}

% --- Marginalien-Inhalt für Recto und Verso ---
% Auf Recto-Seiten (Marginalspalte rechts): linksbündig zum Strich.
% Auf Verso-Seiten (Marginalspalte links):  rechtsbündig zum Strich.
% Das \vspace*{12pt} kompensiert das `before skip` und das `top` der
% tcolorbox, damit das Symbol mit der ersten Zeile der Aufgabenbox fluchtet.
\newcommand{\markercontentRecto}[4]{%
  \vspace*{12pt}\raggedright\color{#4}%
  {\markerfont\fontsize{24}{26}\selectfont #1}\\[0.05em]
  {\small\textsc{#2}}\\
  {\small\textsc{#3}}%
}
\newcommand{\markercontentVerso}[4]{%
  \vspace*{12pt}\raggedleft\color{#4}%
  {\markerfont\fontsize{24}{26}\selectfont #1}\\[0.05em]
  {\small\textsc{#2}}\\
  {\small\textsc{#3}}%
}

% --- Aufgaben-Umgebung ---
% Marginpar wird mit beiden Inhaltsvarianten aufgerufen (recto/verso);
% LaTeX wählt selbst die passende, je nachdem auf welcher Seite die
% Aufgabe beginnt. Der page-aware Strich entsteht im Overlay der
% aufgabebox und wechselt automatisch beim Seitenumbruch mit.
\newenvironment{aufgabe}[4]%
  {% #1 = Symbol, #2 = Label-Zeile-1, #3 = Label-Zeile-2, #4 = Farbe
   \marginpar[\markercontentVerso{#1}{#2}{#3}{#4}]%
              {\markercontentRecto{#1}{#2}{#3}{#4}}%
   \begin{aufgabebox}{#4}\ignorespaces}
  {\end{aufgabebox}}

% --- Komfort-Wrapper für die drei Aufgabentypen ---
% Konvention: der erste Teil der Aufgabe ist der Titel (fett), gefolgt
% vom Aufgabentext. Der Titel wird hier als Argument übergeben.
%
% Wir nutzen {\bfseries ...\par} statt \textbf{...\par}, weil \textbf
% als Inline-Befehl keinen Absatzumbruch (\par) im Argument verträgt
% ("Paragraph ended before \text@command was complete").
\newenvironment{bleistiftuebung}[2]%
  {\begin{aufgabe}{✏}{Bleistift-}{Übung #1}{cBleistift}{\bfseries #2.\par}}
  {\end{aufgabe}}

\newenvironment{pythoneinheit}[2]%
  {\begin{aufgabe}{⌨}{Python-}{Einheit #1}{cPython}{\bfseries #2.\par}}
  {\end{aufgabe}}

\newenvironment{werkzeugcheck}[1]%
  {\begin{aufgabe}{⚙}{Werkzeug-}{Check}{cWerkzeug}{\bfseries #1.\par}}
  {\end{aufgabe}}

% --- Code-Snippet-Einbindung ---
% \pythoncode{tagN/file.py} bindet eine Python-Datei vollständig ein.
% Pfad ist relativ zu latex/code/. Die .py-Dateien sind eigenständig
% lauffähig, sodass Greta sie 1:1 in Thonny verwenden kann.
\newcommand{\pythoncode}[1]{\inputminted{python}{code/#1}}

% \pythoncodeteil{tagN/file.py}{regionname} zeigt nur eine Region
% aus einer Datei. Region-Marker in der Quelle:
%     # region <name>
%     ... Code ...
%     # endregion
% Greta tippt die Datei sukzessive auf, im Buch zeigen wir pro
% Python-Einheit jeweils nur das aktuell relevante Stück.
%
% Implementierung: extract_region.py schneidet die Region zur
% Build-Zeit in eine temporäre Datei, die minted dann einbindet.
% Der Counter sorgt dafür, dass jede Einbindung ihre eigene Datei
% bekommt — nötig, weil minted das Snippet zur Build-Zeit cached.
\newcounter{snippetcntr}
\makeatletter
\newcommand{\pythoncodeteil}[2]{%
  \stepcounter{snippetcntr}%
  \protected@edef\@snippetfile{.snippets/snippet-\arabic{snippetcntr}.py}%
  \immediate\write18{python3 scripts/extract_region.py code/#1 #2 > \@snippetfile}%
  \inputminted{python}{\@snippetfile}%
}
\makeatother

% =============================================================
% Glossar-Einträge.
%
% Eingebunden in der Präambel (vor \begin{document}). Im Text
% verwendet via \gls{schluessel} an der ersten Erwähnung; weitere
% Auftritte können auch \gls{} sein, der Glossar-Eintrag wird dann
% nur noch normal gesetzt.
%
% Schlüssel sind kurz und sprechend; Anzeigename steht in `name=`.
% =============================================================

\newglossaryentry{modulo}{
  name=Modulo,
  description={Operation \(a \bmod n\): der Rest, der bei der Division
    von \(a\) durch \(n\) bleibt. Universalwerkzeug der
    Codierungstheorie}}

\newglossaryentry{paritaet}{
  name=Parität,
  description={Eigenschaft einer Bitfolge, die sich aus der Anzahl der
    Einsen ergibt: \emph{gerade Parität}, wenn die Anzahl gerade
    ist, sonst \emph{ungerade}. Eine Paritätsprüfung erkennt jede
    ungerade Anzahl gekippter Bits}}

\newglossaryentry{paritaetsbit}{
  name=Paritätsbit,
  description={Zusätzliches Bit, das so gewählt wird, dass die
    Gesamt-Parität (meist gerade) zustande kommt. Erkennt
    Einzel-Bitfehler}}

\newglossaryentry{redundanz}{
  name=Redundanz,
  description={Daten-Anteil, der über die reine Information
    hinausgeht und der Erkennung oder Korrektur von
    Übertragungsfehlern dient. Preis für Robustheit}}

\newglossaryentry{ean13}{
  name=EAN-13,
  description={European Article Number, 13-stelliger Strichcode auf
    fast jedem Konsumgut. Letzte Ziffer ist eine Prüfziffer (Modulo
    10 mit Gewichten 1 und 3)}}

\newglossaryentry{isbn10}{
  name=ISBN-10,
  description={Bis 2007 verwendete 10-stellige Buchnummer. Prüfziffer
    Modulo 11; deshalb gelegentlich „X" als Prüfsymbol nötig}}

\newglossaryentry{luhn}{
  name=Luhn-Algorithmus,
  description={Prüfziffer-Verfahren für Kreditkartennummern (Hans Peter
    Luhn, 1954). Verdoppelt jede zweite Ziffer von rechts und summiert
    die Quersummen}}

\newglossaryentry{hammingdistanz}{
  name=Hamming-Distanz,
  description={Anzahl der Stellen, an denen sich zwei gleich lange
    Wörter unterscheiden. Zentrales Maß für Codes}}

\newglossaryentry{mindestdistanz}{
  name=Mindestdistanz,
  description={Kleinste Hamming-Distanz zwischen je zwei Codewörtern
    eines Codes. Bestimmt die Korrektur-/Erkennungs-Fähigkeit:
    \(d-1\) Fehler erkennbar, \(\lfloor (d-1)/2 \rfloor\)
    korrigierbar}}

\newglossaryentry{code}{
  name=Code,
  description={Menge gültiger Codewörter. Die Hamming-Distanzen
    zwischen ihnen bestimmen die Korrektureigenschaften}}

\newglossaryentry{codewort}{
  name=Codewort,
  description={Element eines Codes, also eine gültig codierte Bit-
    oder Symbolfolge}}

\newglossaryentry{hammingschranke}{
  name=Hamming-Schranke,
  description={Untere Grenze für die Anzahl Bits, die zur
    Einzelfehler-Korrektur nötig sind: \(2^k \cdot (n+1) \le 2^n\).
    Wird vom \((7,4)\)-Hamming-Code mit Gleichheit erreicht („perfekter
    Code")}}

\newglossaryentry{hammingcode}{
  name=Hamming-Code,
  description={Familie linearer Codes (Richard Hamming, 1950), die
    einen Einzel-Bitfehler korrigieren. Der \((7,4)\)-Variante codiert
    4 Datenbits in 7 Codebits}}

\newglossaryentry{secded}{
  name=SECDED,
  description={Single Error Correction, Double Error Detection.
    Erweiterung eines Hamming-Codes um ein Gesamt-Paritätsbit:
    erkennt Doppelfehler zusätzlich zur Einzelfehler-Korrektur. In
    ECC-RAM Standard}}

\newglossaryentry{buendelfehler}{
  name=Bündelfehler,
  description={Mehrere benachbarte Bitfehler, etwa durch Kratzer auf
    einer CD oder Funkfading. Realer als zufällige Einzelfehler;
    motiviert Interleaving}}

\newglossaryentry{interleaving}{
  name=Interleaving,
  description={Verschachtelung von Codewörtern vor der Übertragung,
    sodass ein Bündelfehler aus jedem Codewort nur ein einziges Bit
    trifft. Steckt in jeder CD und vielen Funkstandards}}

\newglossaryentry{crc}{
  name=CRC,
  description={Cyclic Redundancy Check. Standardverfahren der
    Fehlererkennung. Daten werden als Polynom interpretiert und
    modulo eines Generatorpolynoms gerechnet}}

\newglossaryentry{generatorpolynom}{
  name=Generatorpolynom,
  description={Festes Polynom \(G(x)\) bei CRC oder Reed-Solomon, das
    die Code-Eigenschaften bestimmt. Codewörter sind durch \(G(x)\)
    teilbar}}

\newglossaryentry{koerper}{
  name=Körper,
  description={Mathematische Struktur, in der man uneingeschränkt
    addieren, subtrahieren, multiplizieren und dividieren kann
    (Ausnahme: Division durch Null). Beispiele: \(\mathbb{R}\),
    \(\mathbb{Q}\), \(\mathbb{Z}/p\mathbb{Z}\) für \(p\) prim}}

\newglossaryentry{endlicherkoerper}{
  name=Endlicher Körper,
  description={Körper mit endlich vielen Elementen. Tritt nur in
    Größen \(p^n\) auf (\(p\) prim). Notation \(GF(p^n)\) für
    \emph{Galois-Feld}. Für Reed-Solomon im Datamatrix wird
    \(GF(2^8)\) verwendet}}

\newglossaryentry{galoisfeld}{
  name=Galois-Feld,
  description={Synonym für endlicher Körper, benannt nach Évariste
    Galois (1811--1832). Geschrieben \(GF(p^n)\)}}

\newglossaryentry{charakteristik}{
  name=Charakteristik,
  description={Kleinste positive Zahl \(p\), für die \(p \cdot 1 = 0\)
    gilt. In jedem \(GF(2^n)\) ist die Charakteristik 2 — daraus
    folgt der „Freshman's Dream" \((a+b)^2 = a^2 + b^2\)}}

\newglossaryentry{irreduziblespolynom}{
  name=Irreduzibles Polynom,
  description={Polynom, das nicht in zwei Polynome niedrigeren Grades
    zerlegt werden kann. Analoges Konzept zur Primzahl. Wird gebraucht,
    um \(GF(2^n)\) als Körper zu bauen}}

\newglossaryentry{inverses}{
  name=Multiplikatives Inverses,
  description={Zu einem Element \(a \ne 0\) eines Körpers das
    Element \(a^{-1}\) mit \(a \cdot a^{-1} = 1\). Existiert in
    Körpern für jedes Element außer Null}}

\newglossaryentry{stuetzstelle}{
  name=Stützstelle,
  description={Punkt \(x_i\), an dem ein Polynom \(f\) ausgewertet
    wird, um den Wert \(f(x_i)\) zu erhalten. Bei Reed-Solomon werden
    \(n\) Stützstellen verwendet, von denen \(k\) zur Rekonstruktion
    reichen}}

\newglossaryentry{reedsolomon}{
  name=Reed-Solomon-Code,
  description={Fehlerkorrigierender Code über \(GF(2^n)\). Daten
    werden als Polynom interpretiert und an \(n\) Stützstellen
    ausgewertet; die \(n\) Werte sind das Codewort. Korrigiert bis
    zu \(\lfloor (n-k)/2 \rfloor\) Symbolfehler}}

\newglossaryentry{aushloeschung}{
  name=Auslöschung,
  description={Fehler mit bekannter Position (engl.\ \emph{erasure}).
    Im Gegensatz zum allgemeinen Fehler („wir wissen nicht, welches
    Symbol verfälscht ist") doppelt so günstig zu korrigieren}}

\newglossaryentry{horner}{
  name=Horner-Schema,
  description={Effizientes Verfahren zur Polynom-Auswertung:
    \(a_0 + x(a_1 + x(a_2 + \dots))\). Spart Multiplikationen
    gegenüber der naiven Auswertung}}


% Defaults, falls ein Master vergisst sie zu setzen.
\providecommand{\maxetappe}{8}
\providecommand{\hauptbuchtitelseite}{}

\begin{document}

\frontmatter
\hauptbuchtitelseite
\pagestyle{plain}

\chapter*{Vorwort}
\addcontentsline{toc}{chapter}{Vorwort}

\section*{Worum es geht}

Auf jeder Schokoladenpackung, in jedem Funkpaket, auf jeder
zerkratzten CD und in jedem Datamatrix-Code auf einem Versandetikett
arbeitet im Hintergrund \emph{Mathematik}, die Fehler erkennt und
korrigiert. Dieses Buch baut diese Mathematik Stück für Stück auf —
von der einfachsten Prüfziffer bis zum vollständigen Datamatrix-Code
mit Reed-Solomon.

\section*{Über Greta}

Das Buch ist als zweiwöchiges Praktikum gedacht und richtet sich
direkt an „Greta", eine Schülerin der zehnten Klasse. Greta gibt es
wirklich; aber selbst wenn du keine Schülerin namens Greta bist, kannst
du das Buch als Greta lesen. Der direkte Ton ist Absicht: niemand
erklärt etwas Schwieriges in der dritten Person.

Vorwissen, das wir voraussetzen: Schulmathematik bis Klasse 10
(Bruchrechnen, Funktionen, ein bisschen Beweise mitdenken) und etwas
Python-Erfahrung (\texttt{for}-Schleifen, Listen, Funktionen — oder die
Bereitschaft, das nebenbei aufzuholen).

\section*{Wie das Buch aufgebaut ist}

Jedes Kapitel ist ein \emph{Tag}. Ein Tag dauert ungefähr zwei Stunden
und besteht aus mehreren Blöcken mit Zeitangaben. In jedem Block
wechseln sich drei Aufgabentypen ab, die du am Strich und Symbol auf
der Seitenaußenkante erkennst:

\begin{itemize}
\item \textcolor{cBleistift}{\textbf{Bleistift-Übungen}} (orange) — mit
  Stift und Karopapier. Selbstdenken vor Selbstrechnen vor Selbstprüfen.
\item \textcolor{cPython}{\textbf{Python-Einheiten}} (blau) — Code in
  Thonny abtippen, ausführen, durchspielen. Computer brute-forcen
  schneller als wir; das ist eines unserer Lieblingsspielzeuge.
\item \textcolor{cWerkzeug}{\textbf{Werkzeug-Checks}} (grün) — kurze
  Setup-Schritte, damit du nicht beim Installieren von Thonny stecken
  bleibst.
\end{itemize}

Im Anhang stehen die \textbf{Lösungen} (Anhang A) und ein
\textbf{Glossar} aller Fachbegriffe. Lösungen sind absichtlich
ausführlich; den Aha-Moment hat man trotzdem nur, wenn man sie erst
nach eigenem Versuch liest.

\section*{Material}

\begin{itemize}
\item Bleistift, Radiergummi, kariertes Papier
\item drei Buntstifte für die Diagramme an Tag 3
\item ein Laptop mit Python (Thonny ist die empfohlene Umgebung)
\item ein Produkt mit Strichcode (Schokolade, Buch — egal) für Tag 1
\end{itemize}

\section*{Tempo}

Die Tagesabschnitte sind als Häppchen gedacht. Wenn ein Tag sich nach
zwei Stunden noch nicht zu Ende anfühlt, ist es kein Verbrechen, ihn
auf zwei Sitzungen aufzuteilen. Tag 5 (endliche Körper) ist
konzeptionell der schwerste — gerade dort darf Pause sein.

Viel Spaß beim Lesen, Rechnen und Tippen.

\tableofcontents

\mainmatter
\pagestyle{scrheadings}

% Etappen 1 .. \maxetappe der Reihe nach einbinden. \include sorgt
% pro Datei für eigenen aux-File und Pagebreak.
\foreach \etappennr in {1,...,\maxetappe}{%
  \include{kapitel/tag\etappennr}%
}

\appendix

% Kapitel-spezifische Anhänge (z. B. Tabellen, Werkstatt-Material).
% Pro Etappe wird anhang/tagN.tex eingebunden, wenn die Datei existiert.
\foreach \etappennr in {1,...,\maxetappe}{%
  \IfFileExists{anhang/tag\etappennr.tex}{%
    \input{anhang/tag\etappennr}%
  }{}%
}

% Lösungen zu allen bisher behandelten Etappen, gesammelt im Anhang.
\chapter{Lösungen}

\begin{quote}
Erst nach eigenem Versuch ansehen — der Aha-Moment beim Selber-Lösen
ist wichtiger als die fertige Lösung am Ende.
\end{quote}

\foreach \etappennr in {1,...,\maxetappe}{%
  \input{loesungen/tag\etappennr}%
}

% \glsaddall zeigt alle Glossar-Einträge (auch die, die im Text nicht
% per \gls{...} referenziert wurden). Solange wir die Markups nicht
% systematisch durch alle Kapitel gezogen haben, ist das die
% praktikable Variante.
\glsaddall
\printglossary[title=Glossar, toctitle=Glossar]

\printindex

\end{document}
